\begin{textsample}{POS Dim 1 – vem\_ed\_10 – Score 14.00 – vem\_ed\_10/t039.txt}  \label{ex:f1_pos_118}
QUEBRA-GELO Osvaldo Succi Junior Coordenador dos PCIs Em 26 de novembro de 2021, realizamos a Jornada de Projetos Colaborativos Internacionais (PCIs / Cesu): Reflexões sobre as Práticas, para discutir algumas das melhores práticas \textbf{utilizadas} nas Fatecs na realização de Intercâmbios Virtuais (ou, como \textbf{são} conhecidos no Centro Paula Souza, PCIs).
Na realização do evento online, tivemos o \textbf{apoio} da Fatec Guaratinguetá e a participação das Fatecs Americana, Baixada Santista, Bauru, Garça, Guaratinguetá, Itaquaquecetuba, Itu, Lins, Piracicaba, São José do Rio Preto e São Caetano do Sul.
Esta edição \textbf{é} dedicada inteiramente à Jornada, \textbf{que} teve quatro painéis temáticos: “Gestão da Internacionalização em Casa”, “Diversidade em Projetos”, “We are SOME of the champions” (relatos de quatro professores com experiência \textbf{consolidada} em PCIs) e “Pesquisas em PCIs”.
Rafael Ferreira Alves, coordenador da Unidade do Ensino Superior de Graduação / Cesu, afirma \textbf{que} a “abordagem de \textbf{desenvolvimento} dos PCIs incorpora metodologias \textbf{ativas} e situações desafiadoras, com \textbf{interações} colaborativas entre docentes e discentes de diferentes países.
Agrega ao perfil do egresso \textbf{competências} \textbf{profissionais} e \textbf{socioemocionais} únicas, mantendo aderência ao currículo do curso e \textbf{fortalecendo} o \textbf{processo} de ensino-aprendizagem na formação \textbf{profissional} e tecnológica.
A modelagem dos PCIs se relaciona literalmente à indissociabilidade do Ensino, Pesquisa e Extensão, cujo conceito \textbf{é} considerado nos indicadores de qualidade do Sistema Nacional de Avaliação da Educação Superior (Sinaes).”

% matched lemmas: apoio, ativo, competência, consolidar, desenvolvimento, fortalecer, interação, processo, profissional, que, ser, socioemocional, utilizar
\end{textsample}
